\section{Evaluation}
\label{sec:evaluation}

The evaluation is a production case study from one organization, not a controlled experiment. The goal is to show what the boundary changed in practice and where the evidence is weaker than we would like.

\subsection{Latency}

Table~\ref{tab:latency} reports \code{/v2/me} full-materialization latency from 1,368 production observations at 512 MB Lambda memory, collected from CloudWatch metrics during March--April 2026. The full path includes JWT verification, membership lookup, tenant record read, impersonation-list assembly, and JSON serialization.

\begin{table}[H]
  \centering
  \small
  \begin{tabularx}{0.65\textwidth}{X r}
    \toprule
    \textbf{Metric} & \textbf{Latency (ms)} \\
    \midrule
    P50 & 856 \\
    P75 & 938 \\
    P90 & 1,038 \\
    P95 & 1,112 \\
    P99 & 1,176 \\
    \bottomrule
  \end{tabularx}
  \caption{Full \code{/v2/me} materialization latency at 512 MB, based on 1,368 production observations from March--April 2026.}
  \label{tab:latency}
\end{table}

ETag-conditional revalidations dominate steady-state traffic at a 91\% hit rate and return \code{304 Not Modified} at 68 ms P50. At that hit rate, amortized per-page cost is approximately 139 ms. Cold starts add 300--600 ms in observed tail cases; Provisioned Concurrency is the obvious mitigation when a stricter latency budget is required \cite{aws_lambda_docs}.

\subsection{Scope-injection checks}

Scope injection is checked both statically and at runtime.

\begin{itemize}
  \item CI verifies that route handler files apply module requirements and that Lambdas receiving \scopeprefix{} reject missing or empty values.
  \item Runtime logs emit actor tenant, effective tenant, and \scopeprefix{} per request. A CloudWatch Insights query surfaces non-impersonation mismatches for review.
\end{itemize}

Within this envelope, structured-log correlation detected zero actor/effective-tenant mismatches across eight months of production operation, from Q3 2025 through April 2026. Monthly impersonation-log audits and customer-reported issues were also clear.

The envelope matters. The method can catch missing scope and actor/effective mismatch. It cannot prove that the control plane never derived a wrong tenant value. If the wrong value is derived consistently, the logs will agree with themselves. We treat that as a limitation, not as an inconvenient detail to hide.

\subsection{Surface filtering}

The platform has 48 documented API endpoints in the measured surface: 5 core endpoints and 43 module-gated endpoints. Filtering the OpenAPI document by enabled modules reduces what limited tenants can discover.

\begin{table}[H]
  \centering
  \small
  \begin{tabularx}{0.72\textwidth}{X r r}
    \toprule
    \textbf{Modules enabled} & \textbf{Endpoints visible} & \textbf{Reduction} \\
    \midrule
    0 (core only) & 5  & 90\% \\
    1             & 16 & 67\% \\
    2             & 27 & 44\% \\
    3+            & 38 & 21\% \\
    Full access   & 48 & 0\% \\
    \bottomrule
  \end{tabularx}
  \caption{Endpoint exposure after OpenAPI filtering. Synthetic module-tier profiles are derived from the 48-endpoint production surface.}
  \label{tab:surface}
\end{table}

Surface filtering does not replace backend authorization. It reduces discovery and support noise. Disabled endpoints are absent from the tenant's spec and return \code{403} if called directly.

\subsection{Governance concentration}

Before the control plane, a governance policy change required deployment to each of 12 data-plane services, averaging roughly four engineer-hours per operation. After the migration, five of six common tenant lifecycle operations became DynamoDB record updates that take under five minutes. Only adding a new module type requires code.

The code concentration is also measurable. Shared JWT validation, role resolution, module gating, and scope derivation live in a 135-line Lambda middleware spine. Reimplementing comparable logic in 12 services would put the same risk in roughly 1,620 lines, before counting drift. The impersonation allow-list check is 67 lines in one place rather than 12 similar versions with slightly different behavior.

Table~\ref{tab:operations} summarizes the operational effects.

\begin{table}[H]
  \centering
  \small
  \begin{tabularx}{\textwidth}{X X X}
    \toprule
    \textbf{Question} & \textbf{Before} & \textbf{After} \\
    \midrule
    Tenant/module changes & Service deployments & Reviewed policy-store updates \\
    Governance logic & Repeated across services & 135-line shared middleware spine \\
    API discovery & Broad generated spec & Filtered by effective modules \\
    Impersonation & Risk of shared or ambiguous context & Actor/effective context logged separately \\
    Migration cost & Ongoing per-service tax & 160 engineer-hours across two sprints \\
    \bottomrule
  \end{tabularx}
  \caption{Operational effects observed after concentrating governance in the control plane.}
  \label{tab:operations}
\end{table}

\subsection{When the pattern pays off}

This pattern is probably not worth the ceremony for a three-service product with simple roles. It starts paying rent when the platform has multiple tenants, multiple modules, support impersonation, and enough services that nobody can confidently name every authorization path from memory. In our deployment, that line was crossed well before 12 services.
